\documentclass[12pt]{article}
\usepackage[T1]{fontenc}
\usepackage[utf8]{inputenc}
\usepackage[spanish,es-nodecimaldot]{babel}
\usepackage{graphicx}
\usepackage[table]{xcolor}
\usepackage{fourier}
\usepackage{multirow,longtable}
\usepackage{stackengine,amsmath,amsfonts,amssymb,braket}
\usepackage{gensymb,wrapfig,enumitem}
\usepackage[pdfborder={0 0 0},colorlinks=true]{hyperref}
\usepackage{fancyhdr}
\pagestyle{fancy}
\usepackage[usegeometry]{typearea}% before geometry!
\usepackage[letterpaper,top=1.16in,left=1.18in,headheight=1in]{geometry}

% --- ABOUT -----
% Formato proyectos para UCM (ver. 1.0)
% -by Javier Rojas (jerojas@ucm.cl)
% Adaptación: Estudio de Caso Unidad 1, INF-321 Computación Numérica

% --- Rellenar datos del encabezado
\title{Cuantización y Propagación del Error en Traducción en Tiempo Real} % Titulo del proyecto
\newcommand{\evaltitle}{ESTUDIO DE CASO --5\%} % Evaluación

\newcommand{\carrera}{Ingeniería Civil Informática} % Carrera
\newcommand{\acurricular}{Computación Numérica (INF-321)} %Actividad Curricular
\newcommand{\academico}{Sergio Hernández} % Academico
\newcommand{\fentrega}{26 de Agosto} % Fecha de Entrega

\newcommand{\pmax}{80} % Puntaje máximo
\newcommand{\pcorte}{48} % Puntaje de Corte

\newcommand{\reglas}{% Inserte las reglas
	\item El trabajo debe presentarse en forma grupal e individual.
	\item Todo experimento debe ser reproducible: indicar hardware (CPU/GPU, sistema operativo),
	versión de \texttt{whisper.cpp} (\textit{commit hash}), versión de Ollama, nombre y \textit{tag} exacto
	de cada modelo, número de hilos, semillas y parámetros de decodificación utilizados. Los resultados
	que no puedan reproducirse no reciben puntaje.
	\item Todo valor numérico reportado debe indicar el número de cifras significativas empleadas y,
	cuando corresponda, si se trata de un error absoluto o relativo. No se aceptan cifras sin unidad
	ni sin definición del error utilizado.
	\item Etiquetar y titular figuras y tablas. Las figuras, por ejemplo, deben llamarse Figura 1, Figura 2,
	etc. (o incluir el número de sección, como Figura 1.3 para la tercera figura de la sección 1) con
	un breve título descriptivo para cada una. Luego, podrás consultarlos fácilmente en el texto.
	Hacer referencia a figuras y tablas por su nombre evita cualquier ambigüedad que algo como
	la figura de arriba pueda introducir.
	\item Las figuras deben estar centradas en la página con un título centrado debajo de ellas. Poner el
	título en cursiva ayuda a que se destaque visualmente del resto del texto. Considere el tamaño
	de cada figura: no deben ser tan pequeñas que los detalles sean difíciles de discernir, ni deben
	ser más grandes de lo necesario para ver claramente todos los detalles, ya que eso desperdicia
	espacio en la página.
	\item Al incluir código o salida de terminal (transcripción de shell) en un documento, formatéelo con
	una fuente monoespaciada. Consolas, DejaVu Sans Mono y Menlo son buenas opciones que se
	encuentran comúnmente disponibles.
	\item No incluya código, graficos, tablas ni una transcripción del shell como captura de pantalla
	del código de una página web, editor de texto o terminal. Hacer esto produce una imagen
	rasterizada del texto que será de baja resolución, pixelada, imposible de resaltar y copiar, e
	inaccesible para cualquiera que utilice un software de lectura de pantalla.
}

\newcommand{\restricciones}{% Inserte restricciones del proyecto
Se aplicará artículo 67º del reglamento del estudiante, el cual indica que, en caso de sorprender copia parcial o exacta, ya sea entre compañeros o reproducidos de algún medio, lo cual implica un 1,0 para todos los involucrados.
Si hay algún requerimiento sin contestar, y en la rúbrica no se puede observar, no aplica puntaje.
Los experimentos deben ejecutarse localmente (\textit{on-device}): no se acepta el uso de servicios de transcripción o traducción en la nube (APIs comerciales) para producir los resultados medidos, ya que el objetivo es controlar la representación numérica de los modelos.
El corpus de audio debe ser de libre uso o de grabación propia con consentimiento de los hablantes; adjuntar la referencia de la fuente.
}

\newcommand{\raprendizaje}{Comprender la naturaleza de los errores numéricos en dispositivos de cálculo, considerando la investigación como estrategia de identificación de problemas.}
\newcommand{\ievaluacion}{
	\begin{itemize}\setlength\itemsep{0pt}
	\item Aplica correctamente la aritmética de punto flotante.
	\item Determina errores de propagación en operaciones.
	\item Maneja dispositivos de cálculo, entendiendo la representación y operatoria de los números.
	\item Utiliza software de apoyo.
	\item Utiliza la investigación en la determinación de problemas de aplicación.
	\end{itemize}
	}

%--- INICIO :: Sección de formato (NO MODIFICAR)

\newcommand*{\useportrait}{%
  \clearpage
  \KOMAoptions{paper=portrait,DIV=current}%switch to portrait
  \newgeometry{% geometry settings for portrait
    letterpaper, left=1.18in, top=1.16in, headheight=1in
  }%
  \fancyhfoffset{0pt}% <- recalculate head and foot width for fancyhdr
}
\newcommand*{\uselandscape}{%
  \clearpage
  \KOMAoptions{paper=landscape,DIV=current}%switch to landscape
  \newgeometry{% geometry settings for landscap
    letterpaper, left=1.18in, top=1.16in, headheight=1in
  }%
  \fancyhfoffset{0pt}% recalculate head and foot width for fancyhdr
}

\fancyhead{}
%\fancyhead[L]{\includegraphics[height=0.6in]{Logo UCM.png}}
\renewcommand{\headrulewidth}{0pt}

\begin{document}
{
  \makeatletter
  \noindent%
  \begin{center}
    {\Large\bf
      \MakeUppercase{\evaltitle}\\%
      ESTUDIO DE CASO\\%
      \MakeUppercase{``\@title''}%
    }
  \end{center}
  \vspace{0.5cm}
  \noindent
  \begin{tabular}{|p{0.3\textwidth}|p{0.69\textwidth}|}
    \hline
    \cellcolor{cyan}\textbf{\color{white}Nombre de los Integrantes}&  \\\hline
    \cellcolor{cyan}\textbf{\color{white}Fecha de Entrega}& \fentrega \\\hline
    \cellcolor{cyan}\textbf{\color{white}Carrera}& \carrera \\\hline
    \cellcolor{cyan}\textbf{\color{white}Académico}& \academico \\\hline
    \cellcolor{cyan}\textbf{\color{white}Actividad Curricular}&\acurricular \\\hline
  \end{tabular}
  \vspace{0.5cm}\\
  \noindent
  \begin{tabular}{|p{0.12\textwidth}|p{0.06\textwidth}|p{0.13\textwidth}|p{0.06\textwidth}|p{0.13\textwidth}|p{0.06\textwidth}|p{0.18\textwidth}|p{0.08\textwidth}|}
    \hline
    \cellcolor{cyan}\textbf{\color{white}PUNTAJE\newline MÁXIMO}&\pmax&\cellcolor{cyan}\textbf{\color{white}PUNTAJE\newline DE CORTE}&\pcorte&\cellcolor{cyan}\textbf{\color{white}PUNTAJE\newline OBTENIDO}&&\cellcolor{cyan}\textbf{\color{white}CALIFICACIÓN}&\\\hline
  \end{tabular}
  \vspace{0.5cm}\\
  \noindent
  \begin{tabular}{|p{0.25\textwidth}|p{0.74\textwidth}|}
    \hline
    \cellcolor{cyan}\textbf{\color{white}Resultados de Aprendizajes Evaluados:}&\raprendizaje\\\hline
    \cellcolor{cyan}\textbf{\color{white}Indicadores de\newline Evaluación:}&\ievaluacion\\\hline
  \end{tabular}\\\vspace{0.1cm}
  \newpage
  \begin{longtable}{|p{1.019\textwidth}|@{}c@{}}
    \hline
    \centering\cellcolor{cyan}\textbf{\color{white}INSTRUCCIONES PARA EL DESRROLLO DE LA EVALUACIÓN}&\\\hline
    {\small
      \vspace{7pt}Para el desarrollo de su evaluación debe considerar lo siguiente:%
      \setlist[enumerate]{itemsep=0pt,parsep=0pt,partopsep=0pt,topsep=0pt}
      \begin{enumerate}
        \reglas
      \end{enumerate}
    }&\\
    {\small
      \setlist[enumerate]{}
      \underline{Restricciones}\newline
      \noindent\restricciones
    }&\\\hline
  \end{longtable}
}

%--- FIN :: Sección de formato (NO MODIFICAR)
\newpage
\noindent
\begin{longtable}{|p{\textwidth}|@{}c@{}}
  \hline
  \centering\cellcolor{cyan}\textbf{\color{white}DESCRIPCIÓN DEL CASO Y DE LAS PREGUNTAS O REQUERIMIENTOS A SER RESUELTOS}&\\\hline
  \endhead
  \textbf{Caso:}&\\
  Un centro de atención de urgencia de la Región del Maule desea implementar un sistema de
  \textbf{traducción en tiempo real} que permita atender a pacientes que no hablan español (por ejemplo,
  creole haitiano o inglés). Por razones de confidencialidad de los datos clínicos, el sistema debe
  ejecutarse \textbf{completamente en el dispositivo local} (un computador de escritorio sin GPU dedicada),
  con una latencia inferior al tiempo real (\textit{real-time factor} $<1$). La arquitectura propuesta es
  una \textbf{cascada} de dos etapas:&\\
  \begin{enumerate}\setlength\itemsep{0pt}
  \item \textbf{Reconocimiento automático del habla} con \texttt{whisper.cpp}
  (\url{https://github.com/ggml-org/whisper.cpp}), implementación en C/C++ del modelo Whisper
  sobre la biblioteca \texttt{ggml}.
  \item \textbf{Traducción y normalización del texto} con un modelo de lenguaje local servido mediante
  \textbf{Ollama} (\url{https://ollama.com}).
  \end{enumerate}&\\


  El equipo disponible no tiene memoria suficiente para ejecutar ambos modelos en precisión completa,
  por lo que se recurre a la \textbf{cuantización}: los pesos, originalmente almacenados en punto flotante
  IEEE-754 de 32 bits (\texttt{binary32}, con $p=24$ bits de mantisa y unidad de redondeo
  $u=2^{-24}\approx 5{,}96\times10^{-8}$), se convierten a formatos de menor ancho de palabra
  (\texttt{F16}, \texttt{Q8\_0}, \texttt{Q5\_1}, \texttt{Q4\_0}, \texttt{Q4\_K\_M}, entre otros).
  En los esquemas \texttt{ggml}/\texttt{gguf} los pesos se agrupan en bloques de $32$ valores que comparten
  un factor de escala $d$ almacenado en punto flotante de 16 bits, de modo que cada peso se reconstruye como&\\
  \vspace{-4pt}
  \[
  \tilde{w}_i = d\,q_i + m, \qquad q_i \in \mathbb{Z}\cap[q_{\min},q_{\max}],
  \]
  \vspace{-14pt}&\\
  lo que introduce un error de representación $\varepsilon_i=\tilde{w}_i-w_i$ acotado por medio paso de
  cuantización, $|\varepsilon_i|\le d/2$. 
  
  &\\
   Una versión preliminar del caso, con instrucciones de instalación y ejecución de los modelos, se encuentra en el repositorio de GitHub:
  \url{https://github.com/shernandez-ucm/estudio_casos_asr}
  &\\\hline
\end{longtable}

\newpage
\noindent
\begin{longtable}{|p{\textwidth}|@{}c@{}}
  \hline
  \centering\cellcolor{cyan}\textbf{\color{white}DESCRIPCIÓN DEL CASO Y DE LAS PREGUNTAS O REQUERIMIENTOS A SER RESUELTOS}&\\\hline
  \endhead
  El centro de salud le solicita a su equipo un informe técnico que responda: \textit{¿cuánta precisión
  numérica se puede sacrificar antes de que la traducción deje de ser clínicamente confiable?}
   &\\
  Para que el equipo pueda cuantificar, y no solo describir, el efecto de estos errores numéricos sobre la
  calidad del resultado final, se emplean tres métricas estándar de evaluación: la \textbf{Tasa de Error de
  Palabra} (\textit{Word Error Rate}, \textbf{WER}) y la \textbf{Tasa de Error de Carácter} (\textit{Character
  Error Rate}, \textbf{CER}), que cuantifican la distancia de edición (inserciones, sustituciones y
  eliminaciones) entre la transcripción generada por \texttt{whisper.cpp} y la transcripción de referencia,
  normalizada por la longitud de esta última; y el \textbf{BLEU} (\textit{Bilingual Evaluation Understudy}),
  que evalúa la calidad de la traducción producida por el modelo servido en Ollama mediante la coincidencia de
  \textit{n}-gramas con una o más traducciones de referencia. Estas métricas permiten trazar, para cada formato
  de cuantización, una curva de degradación de la calidad y así identificar el punto en que el error numérico
  deja de ser despreciable desde el punto de vista clínico.
  &\\
  &\\

  &\\\hline
    \textbf{REQUERIMIENTOS:}
  \vspace{3pt}
  Cada grupo (2 a 3 integrantes) debe diseñar y ejecutar un experimento numérico que cuantifique el
  efecto de la cuantización y la propagación del error en la cascada ASR $\to$ LLM, y comunicar sus
  resultados en un informe y una presentación
  &\\&\\\hline
\end{longtable}

\newpage
\noindent
\begin{longtable}{|p{\textwidth}|@{}c@{}}
  \hline
  \centering\cellcolor{cyan}\textbf{\color{white}REQUERIMIENTOS ESPECÍFICOS}&\\\hline
  \endhead
  \textbf{R1. Representación numérica y aritmética de punto flotante.}&\\
  \begin{enumerate}[label=\alph*)]\setlength\itemsep{0pt}
  \item Compilar \texttt{whisper.cpp} y descargar un modelo (se sugiere \texttt{base} o \texttt{small}) en
  formato \texttt{F32}/\texttt{F16}. Generar al menos tres versiones cuantizadas con la herramienta
  \texttt{quantize} incluida en el repositorio (por ejemplo \texttt{Q8\_0}, \texttt{Q5\_1}, \texttt{Q4\_0}).
  \item Para cada formato, reportar: tamaño en disco, memoria RAM utilizada, bits efectivos por peso
  (incluyendo el costo amortizado del factor de escala por bloque) y el paso de cuantización $d$.
  %\item Deducir analíticamente la cota de error relativo de representación de cada formato y compararla con
  %la unidad de redondeo de \texttt{binary32} y \texttt{binary16}. Verificar experimentalmente el
  %$\varepsilon$ de máquina de ambos formatos en Julia, Python u Octave (\texttt{eps}, \texttt{nextfloat}).
  %\item Extraer una capa de pesos del modelo original y de su versión cuantizada, y calcular
  %empíricamente el error absoluto máximo, el error relativo medio y el error relativo en norma,
  %$\|\mathbf{w}-\tilde{\mathbf{w}}\|_2/\|\mathbf{w}\|_2$. Contrastar con la cota teórica de (c).
  \end{enumerate}&\\
  \textbf{R2. Efecto de la cuantización sobre la precisión de la transcripción y la traducción.}&\\
  \begin{enumerate}[label=\alph*)]\setlength\itemsep{0pt}
  \item Construir un corpus de al menos $10$ segmentos de audio ($15$--$30$ s) con transcripción y
  traducción de referencia (grabación propia o corpus abierto tipo Common Voice o LibriSpeech).
  \item Ejecutar la transcripción y la traducción con cada formato del modelo, midiendo
  \textbf{WER} y \textbf{CER} para la transcripción, y \textbf{BLEU} o \textbf{chrF} para la traducción,
  además del \textit{real-time factor} y el consumo de memoria.
  \item Graficar la curva de precisión \textit{versus} bits por peso e identificar el punto de quiebre a
  partir del cual el error crece de manera no lineal. Discutir el compromiso
  precisión--memoria--latencia usando lenguaje de análisis de error (error absoluto y relativo, error
  de truncamiento frente a error de redondeo).
  \end{enumerate}&\\
  \textbf{R3. Propagación del error en la cascada con Ollama.}&\\
  \begin{enumerate}[label=\alph*)]\setlength\itemsep{0pt}
  \item Servir un mismo modelo de lenguaje en Ollama con al menos dos cuantizaciones distintas
  (por ejemplo \texttt{q4\_0} y \texttt{q8\_0} de un modelo de 7--8B) y usarlo para traducir, fijando
  \texttt{temperature} $=0$ para aislar el efecto numérico del efecto del muestreo.
  \item Separar las fuentes de error mediante tres condiciones experimentales:
  (i) traducir la \textbf{transcripción de referencia} (error del LLM aislado);
  (ii) traducir la salida del ASR en alta precisión; (iii) traducir la salida del ASR cuantizado
  (error compuesto). Reportar los tres errores en una misma tabla.
  %\item Estimar un \textbf{factor de amplificación} empírico
  %$\kappa \approx (\text{error relativo de salida})/(\text{error relativo de entrada})$
  %y discutir si la cascada amplifica, atenúa o mantiene el error. Relacionar con el concepto de número
  %de condición y de problema bien o mal condicionado.
  %\item Repetir al menos un experimento variando el número de hilos (\texttt{-t}) y la semilla, dejando fija
  %la entrada, para evidenciar el efecto de la \textbf{no asociatividad} de la suma en punto flotante sobre
  %la reproducibilidad de los resultados.
  \end{enumerate}&\\
  \textbf{R4. Investigación e identificación del problema.}&\\
  \begin{enumerate}[label=\alph*)]\setlength\itemsep{0pt}
 % \item Cada integrante debe seleccionar y revisar críticamente \textbf{un} artículo o documentación
 % técnica sobre cuantización, aritmética de baja precisión o estabilidad numérica en inferencia
 % (por ejemplo, sobre formatos \texttt{FP16}, \texttt{BF16} o \texttt{INT8}, esquemas \textit{k-quants},
 % o \textit{outliers} de activación), indicando qué aporta al caso analizado.
  \item Formular una recomendación técnica al centro de salud: qué formato usar, con qué margen de
  error esperado y qué riesgos residuales existen para un uso clínico.
  \end{enumerate}&\\&\\\hline
\end{longtable}

\newpage
\noindent
\begin{longtable}{|p{\textwidth}|@{}c@{}}
	\hline
	\centering\cellcolor{cyan}\textbf{\color{white}FORMATO DE ENTREGA}&\\\hline
	\endhead
	\textbf{Formato:}&\\
	La entrega consta de un \textbf{informe breve} (máximo 6 páginas, sin contar anexos) y una
	\textbf{presentación grupal de 10 minutos}, con la siguiente estructura:&\\
	\begin{itemize}\setlength\itemsep{0pt}
		\item Introducción y planteamiento del problema numérico (Grupal).
		\item Marco teórico: representación en punto flotante, cuantización y propagación del error (Grupal).
		\item Diseño experimental, software utilizado y resultados en tablas y figuras (Grupal).
		\item Análisis de la propagación del error en la cascada y conclusiones (Grupal).
		\item Revisión crítica de un artículo o documentación técnica (Individual).
	\end{itemize}&\\
	Una vez finalizada la presentación grupal, cada integrante debe presentar su artículo (1 diapositiva)
	demostrando capacidad de análisis, síntesis, evaluación y aplicación del conocimiento.
	En total cada grupo dispone de 20 minutos para la presentación grupal e individual.&\\
	Se debe adjuntar además un repositorio o carpeta comprimida con los \textit{scripts} utilizados
	(Julia, Python u Octave), los archivos de audio o su referencia, y las salidas crudas de cada corrida.
	Se sugiere organizar los resultados en tablas del siguiente tipo:&\\
	\vspace{6pt}\small
	\hspace*{\fill}\begin{tabular}{|p{0.17\textwidth}|p{0.12\textwidth}|p{0.12\textwidth}|p{0.09\textwidth}|p{0.09\textwidth}|p{0.09\textwidth}|}
	\hline
	\textbf{Formato} & \textbf{Bits/peso} & \textbf{Memoria (MB)} & \textbf{WER (\%)} & \textbf{chrF} & \textbf{RTF} \\\hline
	F32 & 32 & & & & \\\hline
	F16 & 16 & & & & \\\hline
	Q8\_0 & & & & & \\\hline
	Q5\_1 & & & & & \\\hline
	Q4\_0 & & & & & \\\hline
	\end{tabular}\hspace*{\fill}\vspace{6pt}&\\
	\vspace{6pt}\small
	\hspace*{\fill}\begin{tabular}{|p{0.32\textwidth}|p{0.15\textwidth}|p{0.15\textwidth}|p{0.15\textwidth}|}
	\hline
	\textbf{Condición experimental} & \textbf{Error de entrada} & \textbf{Error de salida} & \textbf{$\kappa$ empírico} \\\hline
	Referencia $\to$ LLM & $0$ & & \\\hline
	ASR alta precisión $\to$ LLM & & & \\\hline
	ASR cuantizado $\to$ LLM & & & \\\hline
	\end{tabular}\hspace*{\fill}\vspace{6pt}&\\&\\\hline
\end{longtable}

\uselandscape
\section*{RÚBRICA}
%\noindent
\begingroup
\scriptsize
\setlength{\LTpre}{0pt}\setlength{\LTpost}{0pt}
\begin{longtable}{|p{0.15\textwidth}|p{0.15\textwidth}|p{0.15\textwidth}|p{0.15\textwidth}|p{0.15\textwidth}|p{0.13\textwidth}|}
  \cline{2-5}
  \multicolumn{1}{p{0.15\textwidth}}{}&\multicolumn{4}{|c|}{Valores}&\multicolumn{1}{p{0.13\textwidth}}{}\\\hline
  Indicadores&\centering Destacado\newline (10 pts)&\centering Competente\newline (7 pts)&\centering Básico\newline (5 pts)&\centering En Desarrollo\newline (2 pts)&Observaciones\\\hline
  \endhead

Presentación y formato (Grupal) &Informe y presentación impecables, bien organizados, dentro del tiempo, con tablas y figuras rotuladas y sin errores tipográficos.&Buena presentación, mayormente organizada, con pocos errores tipográficos o de rotulación.&Presentación aceptable, pero con errores tipográficos o de organización.&Presentación desorganizada o con numerosos errores.&1 \vspace{0.15cm}~\\\hline

Representación numérica y punto flotante (Grupal) &Deduce correctamente las cotas de error de cada formato, las verifica experimentalmente y las contrasta con la unidad de redondeo de \texttt{binary32}/\texttt{binary16}.&Deduce y verifica las cotas de error de los formatos, con contraste parcial.&Describe los formatos y reporta tamaños, pero sin deducir ni verificar cotas de error.&No distingue los formatos ni aplica la aritmética de punto flotante.&2 \vspace{0.15cm}~\\\hline

Diseño experimental y uso de software (Grupal)&Experimento reproducible y bien controlado; usa \texttt{whisper.cpp} y Ollama con métricas apropiadas, corpus adecuado y \textit{scripts} documentados.&Experimento correcto y reproducible, con métricas apropiadas pero controles o documentación parciales.&Ejecuta las herramientas y reporta resultados, sin control de variables ni reproducibilidad.&No logra ejecutar los experimentos o los resultados no son verificables.&3 \vspace{0.15cm}~\\\hline

Manejo de dispositivos de cálculo (Grupal)&Reporta memoria, \textit{real-time factor} y bits por peso de cada formato, y evidencia el efecto del número de hilos y del orden de las reducciones sobre la salida.&Reporta memoria y latencia por formato y verifica la variabilidad al cambiar hilos o semilla.&Reporta memoria o latencia sin analizar su relación con la representación numérica.&No caracteriza el dispositivo ni el costo de cada formato.&4 \vspace{0.15cm}~\\\hline

Análisis de propagación del error (Grupal)&Separa las fuentes de error, estima el factor de amplificación, lo interpreta en términos de condicionamiento y contrasta con las cotas teóricas.&Separa las fuentes de error y estima la amplificación, con interpretación parcial.&Compara errores entre formatos sin descomponer ni cuantificar la propagación.&Análisis del error ausente o incorrecto.&5 \vspace{0.15cm}~\\\hline

Investigación y evaluación crítica (Individual)&Selecciona una fuente pertinente, evalúa críticamente sus fortalezas y debilidades y la vincula explícitamente con los resultados obtenidos.&Evalúa críticamente la fuente y la relaciona con el caso.&Resume la fuente e identifica alguna fortaleza o debilidad.&Revisión ausente, superficial o sin relación con el caso.&6 \vspace{0.15cm}~\\\hline

Claridad y coherencia (Individual)&Presenta las ideas de manera clara, concisa y lógica, con lenguaje técnico apropiado y uso correcto de cifras significativas.&Ideas claras y bien organizadas, con lenguaje técnico apropiado.&Utiliza un lenguaje apropiado, con exposición poco estructurada.&Exposición confusa o insuficiente.&7 \vspace{0.15cm}~\\\hline

Conclusiones y recomendación técnica (Individual)&Conclusión fundamentada en la evidencia numérica, con recomendación explícita, margen de error esperado y riesgos residuales del uso clínico.&Conclusión clara y bien relacionada con los objetivos, con recomendación fundamentada.&Conclusión presente pero poco clara o insuficientemente relacionada con los resultados.&Conclusión ausente o irrelevante.&8 \vspace{0.15cm}~\\\hline

\multicolumn{5}{|r|}{\textbf{PUNTAJE TOTAL (máximo \pmax\ puntos)}}&\vspace{0.15cm}~\\\hline

\end{longtable}
\endgroup
\useportrait
\section*{ESCALA DE NOTAS}
\noindent
Puntaje máximo: \pmax\ puntos. Puntaje de corte (nota 4,0): \pcorte\ puntos, correspondiente a un 60\% de exigencia.
\vspace{0.3cm}
\begin{center}
\small
\begin{tabular}{|p{1.3cm}|p{1.3cm}|p{1.3cm}|p{1.3cm}|p{1.3cm}|p{1.3cm}|p{1.3cm}|p{1.3cm}| }
\hline
 Puntaje&  Nota&  Puntaje&  Nota&  Puntaje&  Nota&  Puntaje&  Nota \\ \hline
0	&	\color{red}1,0	&	21	&	\color{red}2,3	&	42	&	\color{red}3,6	&	63	&	5,4\\ \hline
1	&	\color{red}1,1	&	22	&	\color{red}2,4	&	43	&	\color{red}3,7	&	64	&	5,5\\ \hline
2	&	\color{red}1,1	&	23	&	\color{red}2,4	&	44	&	\color{red}3,8	&	65	&	5,6\\ \hline
3	&	\color{red}1,2	&	24	&	\color{red}2,5	&	45	&	\color{red}3,8	&	66	&	5,7\\ \hline
4	&	\color{red}1,3	&	25	&	\color{red}2,6	&	46	&	\color{red}3,9	&	67	&	5,8\\ \hline
5	&	\color{red}1,3	&	26	&	\color{red}2,6	&	47	&	\color{red}3,9	&	68	&	5,9\\ \hline
6	&	\color{red}1,4	&	27	&	\color{red}2,7	&	48	&	4,0	&	69	&	6,0\\ \hline
7	&	\color{red}1,4	&	28	&	\color{red}2,8	&	49	&	4,1	&	70	&	6,1\\ \hline
8	&	\color{red}1,5	&	29	&	\color{red}2,8	&	50	&	4,2	&	71	&	6,2\\ \hline
9	&	\color{red}1,6	&	30	&	\color{red}2,9	&	51	&	4,3	&	72	&	6,3\\ \hline
10	&	\color{red}1,6	&	31	&	\color{red}2,9	&	52	&	4,4	&	73	&	6,3\\ \hline
11	&	\color{red}1,7	&	32	&	\color{red}3,0	&	53	&	4,5	&	74	&	6,4\\ \hline
12	&	\color{red}1,8	&	33	&	\color{red}3,1	&	54	&	4,6	&	75	&	6,5\\ \hline
13	&	\color{red}1,8	&	34	&	\color{red}3,1	&	55	&	4,7	&	76	&	6,6\\ \hline
14	&	\color{red}1,9	&	35	&	\color{red}3,2	&	56	&	4,8	&	77	&	6,7\\ \hline
15	&	\color{red}1,9	&	36	&	\color{red}3,3	&	57	&	4,8	&	78	&	6,8\\ \hline
16	&	\color{red}2,0	&	37	&	\color{red}3,3	&	58	&	4,9	&	79	&	6,9\\ \hline
17	&	\color{red}2,1	&	38	&	\color{red}3,4	&	59	&	5,0	&	80	&	7,0\\ \hline
18	&	\color{red}2,1	&	39	&	\color{red}3,4	&	60	&	5,1	&		&	\\ \hline
19	&	\color{red}2,2	&	40	&	\color{red}3,5	&	61	&	5,2	&		&	\\ \hline
20	&	\color{red}2,3	&	41	&	\color{red}3,6	&	62	&	5,3	&		&	\\ \hline
\end{tabular}
\end{center}
\end{document}
